% stub

自 C++11 引入移动语义以来，传递字符串变得快了许多。然而你仍可能遇到大量临时的字符串拷贝。在 C++17 中，你得到了一个名为 \texttt{string\_view} 的新类型。它允许你创建一个连续字符序列的、常量、非拥有的视图。你可以操作该视图并到处传递它，而无需拷贝被引用的数据。不过，这一特性也带来了一些代价：你需要小心避免产生“悬垂”的视图，而且通常这类视图可能并非以空字符结尾。

本章你将学到：

\begin{itemize}
\item 什么是 \texttt{string\_view}？
\item 它为何能加速你的代码？
\item 使用 \texttt{string\_view} 对象有哪些风险？
\item 什么是引用生命期延长，它对 \texttt{string\_view} 意味着什么？
\item 如何借助 \texttt{string\_view} 让你的 API 更具泛型能力？
\end{itemize}


\section{基础}

我们先做一个小实验：

下面的例子会创建多少个字符串拷贝？

\begin{codebox}[cpp]
// string function:
std::string StartFromWordStr(const std::string& strArg, const std::string& word)
{
return strArg.substr(strArg.find(word)); // substr creates a new string
}
// call:
std::string str {"Hello Amazing Programming Environment" };
auto subStr = StartFromWordStr(str, "Programming Environment");
std::cout << subStr << '\n';
\end{codebox}

你能把它们都数出来吗？

答案是 3 个或 5 个，取决于编译器；但通常应该是 3 个。

\begin{itemize}
\item 第一个是为 \texttt{str} 创建的。
\item 第二个是 \texttt{StartFromWordStr} 的第二个参数——该参数类型是 \texttt{const string\&}，由于我们传入的是 \texttt{const char*}，它会创建一个新的字符串。
\item 第三个来自 \texttt{substr}，它返回一个新的字符串。
\item 此外，当对象从函数返回时，可能还会再多出一两个拷贝。
\end{itemize}

不过通常编译器能够优化并消除这些拷贝（尤其是自 C++17 起，那种情况下拷贝消除已成为强制要求）。

\begin{itemize}
\item 如果字符串较短，那么由于小字符串优化¹，可能根本不会发生堆分配。
\end{itemize}

上面的例子是经过简化的。然而你可以想象一下生产环境中的代码：字符串操作频繁发生，那种场景下，编译器创建的所有临时对象连数都很难数清。

解决临时拷贝问题的一个更好的做法是使用 \texttt{std::string\_view}。顾名思义，你拿到的不再是原始字符串本身，而只是它的一个非拥有视图。大多数情况下，它只是一个指向连续字符序列的指针加上一个长度。你可以到处传递它，并使用大多数常见的字符串操作。

视图与 \texttt{substr} 这类字符串操作配合得很好。在典型情况下，每一次子串操作都会创建另一个更小的字符串拷贝；而使用 \texttt{string\_view} 时，\texttt{substr} 只是映射到原始缓冲区的不同部分，既不额外消耗内存，也不做动态分配。

下面是我们代码的改写版本，它接受 \texttt{string\_view}：

\footnotetext{¹ 小字符串优化（Small String Optimisation）并未在 C++ 标准中规定，但它是主流编译器普遍采用的优化。目前，MSVC（VS 2017）/GCC（8.1）下为 15 个字符，Clang（6.0）下为 22 个字符。}

\begin{codebox}[cpp]
// Chapter string_view/avoiding_copies_string_view.cpp
std::string_view StartFromWord(std::string_view str, std::string_view word)
{
return str.substr(str.find(word)); // substr creates only a new view
}
// call:
std::string str {"Hello Amazing Programming Environment"};
auto subView = StartFromWord(str, "Programming Environment");
std::cout << subView << '\n';
\end{codebox}

在上面的例子中，我们只有一次分配——只为那个主字符串 \texttt{str}。没有任何 \texttt{string\_view} 操作会为新的字符串触发拷贝或额外的内存分配。当然，\texttt{string\_view} 本身会被拷贝——但由于它只是一个指针加一个长度，这比拷贝整个字符串要高效得多。

有一点警告：虽然这个例子展示了字符串视图的优化能力，但请继续往下读，看看这段代码所暗藏的风险与前提假设！或者你现在就已经能发现其中几个问题了？

\subsection{何时使用}

那么，你应当在什么时候使用 \texttt{string\_view}：

\begin{itemize}
\item 优化：你可以仔细审视自己的代码，把各种字符串操作替换成 \texttt{string\_view}。在大多数情况下，你最终会得到更快的代码以及更少的内存分配。
\item 作为 \texttt{const std::string\&} 参数的潜在替代品——尤其是那些不需要拥有权、也不会保存该字符串的函数中。
\item 处理来自其它 API 的字符串：\texttt{QString}、\texttt{CString}、\texttt{const char*}……只要内容是放在一段连续内存中、且具备基本的字符类型，都可以。你可以写一个接受 \texttt{string\_view} 的函数，而从那种其它实现转换过来时不会发生任何转换开销。
\end{itemize}

无论如何，有一点很重要：要记住它仅仅是一个非拥有视图，因此一旦原始对象不复存在，视图就变成了垃圾，你可能会陷入麻烦。

此外，\texttt{string\_view} 可能不包含空终止符，所以你的代码也必须能应对这一点。例如，把一个 \texttt{string\_view} 传给一个接受以空字符结尾的字符串的函数，从来都不是个好主意。关于这一点，我们会在单独的一节——“使用 \texttt{string\_view} 的风险”——中详谈。

\section{\texttt{std::basic\_string\_view} 类型}

虽然我们谈论的是 \texttt{string\_view}，但有一点很重要：它其实只是一个模板类 \texttt{basic\_string\_view} 的特化：

\begin{codebox}[cpp]
template<
class CharT,
class Traits = std::char_traits<CharT>
> class basic_string_view;
\end{codebox}

\texttt{Traits} 类用于抽象出针对字符类型的操作，例如如何比较字符、如何在一个序列中查找某个字符。

这样的层级关系与 \texttt{std::string} 的情况类似——\texttt{std::string} 也是 \texttt{std::basic\_string} 的一个特化。

我们有以下特化版本：

\begin{codebox}[cpp]
std::string_view
std::basic_string_view<char>
std::wstring_view
std::basic_string_view<wchar_t>
std::u16string_view
std::basic_string_view<char16_t>
std::u32string_view
std::basic_string_view<char32_t>
\end{codebox}

如你所见，各特化版本使用了不同的底层字符类型。

为了方便起见，本章余下部分只考虑 \texttt{string\_view}。

\section{\texttt{std::string\_view} 的创建}

你可以通过几种方式创建一个 \texttt{string\_view}：

\begin{itemize}
\item 从 \texttt{const char*}——提供一个指向以空字符结尾的字符串的指针；
\item 从 \texttt{const char*} 加长度；
\item 通过从 \texttt{std::string} 转换得到；
\item 通过使用 \texttt{""sv} 字面量。
\end{itemize}

下面是各种创建方式的示例：

\begin{codebox}[cpp]
// Chapter string_view/string_view_creation.cpp
const char* cstr = "Hello World";
// the whole string:
std::string_view sv1 { cstr };
std::cout << sv1 << ", len: " << sv1.size() << '\n';
// slice
std::string_view sv2 { cstr, 5 }; // not null-terminated!
std::cout << sv2 << ", len: " << sv2.size() << '\n';
// from string:
std::string str = "Hello String";
std::string_view sv3 = str;
std::cout << sv3 << ", len: " << sv3.size() << '\n';
// ""sv literal
using namespace std::literals;
std::string_view sv4 = "Hello\0 Super World"sv;
std::cout << sv4 << ", len: " << sv4.size() << '\n';
std::cout << sv4.data() << " - till zero\n";
\end{codebox}

程序将输出：

\begin{codebox}
Hello World, len: 11
Hello, len: 5
Hello String, len: 12
Hello Super World, len: 18
Hello - till zero
\end{codebox}

请注意最后两行：\texttt{sv4} 在中间包含了一个 \texttt{'\textbackslash0'}，但 \texttt{std::cout} 仍然能把整个序列打印出来。在最后一行中，我们尝试用 \texttt{.data()} 来打印，而得到的是一个字符串指针，于是打印在遇到空终止符处就中断了。

\section{其它操作}

\texttt{string\_view} 被设计得与 \texttt{std::string} 非常相似。不过，由于视图是非拥有的，任何会修改数据的操作都不可能进入它的 API。下面是你可以配合这个新类型使用的方法的简要清单：

\textbf{迭代器：}

\begin{itemize}
\item \texttt{cbegin()}/\texttt{begin()} / \texttt{crbegin()}/\texttt{rbegin()} —— 两者都是 const 且 constexpr；
\item \texttt{cend()}/\texttt{end()} / \texttt{crend()}/\texttt{rend()} —— 两者都是 const 且 constexpr。
\end{itemize}

\textbf{访问元素：}

\begin{itemize}
\item \texttt{operator[]}
\item \texttt{at}
\item \texttt{front}
\item \texttt{back}
\item \texttt{data}
\end{itemize}

\textbf{大小与容量：}

\begin{itemize}
\item \texttt{size}/\texttt{length}
\item \texttt{max\_size}
\item \texttt{empty}
\end{itemize}

\textbf{修改器：}

\begin{itemize}
\item \texttt{remove\_prefix}
\item \texttt{remove\_suffix}
\item \texttt{swap}
\end{itemize}

\textbf{其它：}

\begin{itemize}
\item \texttt{copy}（不是 constexpr）
\item \texttt{substr} —— 复杂度为 O(1)，而不像 \texttt{std::string} 中那样是 O(n)
\item \texttt{compare}
\item \texttt{find}
\item \texttt{rfind}
\item \texttt{find\_first\_of}
\item \texttt{find\_last\_of}
\item \texttt{find\_first\_not\_of}
\item \texttt{find\_last\_not\_of}
\item 用于字典序比较的运算符：\texttt{==}、\texttt{!=}、\texttt{<=}、\texttt{>=}、\texttt{<}、\texttt{>}
\item 用于 \texttt{ostream} 输出的 \texttt{operator<<}
\item 针对 \texttt{std::hash} 的特化
\end{itemize}

关于上述操作，有几件关键的事：

\begin{itemize}
\item 上述所有方法（除 \texttt{copy}、\texttt{operator<<} 以及 \texttt{std::hash} 的特化之外）也都被标记为 constexpr！借助这一能力，你现在可以在常量表达式中处理连续的字符序列了。
\item 上面的清单几乎与所有不可变的字符串操作一致。不过有两个新方法：\texttt{remove\_prefix} 和 \texttt{remove\_suffix}——它们不是 const 的，会修改 \texttt{string\_view} 对象本身。注意，它们仍然无法修改被引用的数据。
\item \texttt{operator[]}、\texttt{at}、\texttt{front}、\texttt{back}、\texttt{data} 也都是 const 的——因此你无法更改底层字符序列（它只是“只读访问”）。在 \texttt{std::string} 中，这些方法有返回引用的重载版本——于是你获得的是“写访问”。这一点在 \texttt{string\_view} 上是不可能的。
\end{itemize}

\begin{myInfoBox}
More in C++20!

关于 C++20，我们至少还会得到两个新方法：

\begin{itemize}
\item \texttt{starts\_with}
\item \texttt{ends\_with}
\end{itemize}

它们同时针对 \texttt{std::string\_view} 和 \texttt{std::string} 实现。截至此刻（2018 年 7 月），Clang 6.0 已经支持这些函数。如果你愿意，可以尝试使用它们。
\end{myInfoBox}

\section{使用 string\_view 的风险}

\texttt{std::string\_view} 加入标准，主要是为了允许进行性能优化。然而，它并不是字符串的替代品！正因如此，当你使用视图时，必须牢记几件潜在的有风险的事：

\paragraph{留意非以空字符结尾的字符串}

\texttt{string\_view} 在字符串末尾可能并不包含 \texttt{\textbackslash0}。所以你必须为此做好准备。

\begin{itemize}
\item \texttt{string\_view} 与所有接受传统 C 风格字符串的函数配合起来都很成问题，因为 \texttt{string\_view} 打破了 C 风格字符串关于终止符的假设。如果一个函数只接受 \texttt{const char*} 参数，那么把 \texttt{string\_view} 传进去很可能是个坏主意。另一方面，当这样的函数同时接受 \texttt{const char*} 和长度参数时，这样做也许是安全的。
\item 转换成字符串——你不仅需要指定指向连续字符序列的指针，还需要指定长度。
\end{itemize}

\paragraph{引用与临时对象}

\texttt{string\_view} 并不拥有内存，因此在处理临时对象时必须非常小心。

一般而言，一个 \texttt{string\_view} 的生命期绝不能超过那个拥有字符串的对象的生命期。在以下情形中这一点尤为重要：

\begin{itemize}
\item 从函数返回 \texttt{string\_view}；
\item 把 \texttt{string\_view} 存进对象或容器中。
\end{itemize}

为了探究所有这些问题，让我们从本章开头的初始示例说起。

\subsection{初始示例的问题}

引言一节向你展示过这样一个示例：

\begin{codebox}[cpp]
std::string_view StartFromWord(std::string_view str, std::string_view word)
{
return str.substr(str.find(word)); // substr creates only a new view
}
\end{codebox}

这段代码对于非以空字符结尾的字符串没有任何问题——因为所有函数都来自 \texttt{string\_view} 的 API。

然而，临时对象又如何呢？

如果你这样调用会发生什么？

\begin{codebox}[cpp]
auto str = "My Super"s;
auto sv = StartFromWord(str + " String", "Super");
// use `sv` later in the code...
\end{codebox}

这样的代码可能会直接崩掉！

\texttt{"Super"} 是一个临时的 \texttt{const char*} 字面量，它被作为 \texttt{string\_view word} 传入函数。这没问题，因为该临时对象保证会存活得与整个函数调用一样久。然而，字符串拼接 \texttt{str + " String"} 的结果是一个临时对象，而函数却把这个临时对象的 \texttt{string\_view} 返回到了调用之外！

因此，这类情况的一般性建议是：虽然从一个函数返回 \texttt{string\_view} 是可能的，但你必须小心，并且要确切清楚底层字符串的状态。

要理解临时值的问题，最好看一看引用生命期延长。

\subsection{引用生命期延长}

下面这种情况会发生什么？

\begin{codebox}[cpp]
std::vector<int> GenerateVec()
{
return std::vector<int>(5, 1);
}
const std::vector<int>& refv = GenerateVec();
\end{codebox}

上面的代码安全吗？

是的——C++ 规则规定，绑定到一个 const 引用的临时对象的生命期会被延长到与该引用本身一样久。

下面是从标准中摘录的一个完整示例（C++17 草案 N4687²）15.2 临时对象 [class.temporary]：

\footnotetext{² \url{https://wg21.link/n4687}}

\begin{myQuote}
[Example:
struct S {
S();
S(int);
friend S operator+(const S\&, const S\&);
~S();
};
S obj1;
const S\& cr = S(16)+S(23);
S obj2;
The expression S(16) + S(23) creates three temporaries: a first temporary T1 to hold the result of
the expression S(16), a second temporary T2 to hold the result of the expression S(23), and a third
temporary T3 to hold the result of the addition of these two expressions. The temporary T3 is then
bound to the reference cr. It is unspecified whether T1 or T2 is created first. On an implementation
where T1 is created before T2, T2 shall be destroyed before T1. The temporaries T1 and T2 are
bound to the reference parameters of operator+; these temporaries are destroyed at the end of the
full-expression containing the call to operator+. The temporary T3 bound to the reference cr is
destroyed at the end of cr’s lifetime, that is, at the end of the program. In addition, the order in
which T3 is destroyed takes into account the destruction order of other objects with static storage
duration. That is, because obj1 is constructed before T3, and T3 is constructed before obj2, obj2
shall be destroyed before T3, and T3 shall be destroyed before obj1. -end example]
\end{myQuote}

虽然你最好不要为所有变量都写出这样的代码，但在某些情况下它确实是个便利的特性，例如：

\begin{codebox}[cpp]
for (auto &elem : GenerateVec())
{
// ...
}
\end{codebox}

在上面这个例子中，\texttt{GenerateVec} 在基于范围的 for 循环内部被绑定到一个引用（用于 vector 起始处的右值引用）。如果没有生命期延长这一支持，这段代码就会崩溃。

它与 \texttt{string\_view} 有何关联？

对于 \texttt{string\_view} 而言，下面这段代码通常很容易出错：

\begin{codebox}[cpp]
std::string func()
{
std::string s;
// build s...
return s;
}
std::string_view sv = func();
// no temp lifetime extension!
\end{codebox}

这可能并不显而易见——\texttt{string\_view} 同样是一个常量视图，所以本应表现得几乎像一个 const 引用。但根据现有的 C++ 规则，事实并非如此——编译器会在整个表达式结束后立即销毁那个临时对象。在这种情况下，生命期无法被延长。

\texttt{string\_view} 只是一个代理对象，类似于另一段代码：

\begin{codebox}[cpp]
std::vector<int> CreateVector() { ... }
std::string GetString() { return "Hello"; }
auto &x = CreateVector()[10]; // arbitrary element!
auto pStr = GetString().c_str();
\end{codebox}

在这两种情况下，\texttt{x} 和 \texttt{pStr} 都不会延长在 \texttt{CreateVector()} 或 \texttt{GetString()} 中创建的临时对象的生命期。

你可以通过以下方式修复它：

\begin{codebox}[cpp]
std::string func()
{
std::string s;
// build s...
return s;
}
auto temp = func();
std::string_view sv { temp };
// fine lifetime of temporary is extended through `temp`
\end{codebox}

每当你把某个函数的返回值赋给别的东西时，都必须确保该对象的生命期是正确的。

有一份提案致力于修复 \texttt{string\_view} 以及其它本应具有引用生命期延长语义的类型所存在的问题：参见 P0936³。

\footnotetext{³ \url{https://wg21.link/p0936}}

\section{从 string\_view 初始化字符串成员}

由于 \texttt{string\_view} 在作为函数传入参数时是 \texttt{const string\&} 的潜在替代品，我们不妨考虑字符串成员初始化的情况。在这里，\texttt{string\_view} 是最佳候选吗？

例如：

\begin{codebox}[cpp]
class UserName {
std::string mName;
public:
UserName(const std::string& str) : mName(str) { }
};
\end{codebox}

如你所见，构造函数接受的是 \texttt{const std::string\& str}。

你或许可以把这个常量引用替换成 \texttt{string\_view}：

\begin{codebox}[cpp]
UserName(std::string_view sv) : mName(sv) { }
\end{codebox}

我们来比较一下这两种实现在三种情形下的表现：从字符串字面量创建、从左值创建、从右值引用创建：

\begin{codebox}[cpp]
// creation from a string literal
UserName u1{"John With Very Long Name"};
// creation from l-value:
std::string s2 {"Marc With Very Long Name"};
UserName u2 { s2 };
// use s2 later...
// from r-value reference
std::string s3 {"Marc With Very Long Name"};
UserName u3 { std::move(s3) };
// third case is also similar to taking a return value:
std::string GetString() { return "some string..."; }
UserName u4 { GetString() };
\end{codebox}

现在我们来分析 \texttt{UserName} 构造函数的两个版本——使用字符串引用，或使用字符串视图。

请注意，\texttt{s2} 和 \texttt{s3} 的分配/创建不计入考量，我们只关注构造函数调用时发生了什么。对于 \texttt{s2}，我们也可以假设它之后还会在代码中被使用。

对于 \texttt{const std::string\&}：

\begin{itemize}
\item \texttt{u1}——两次分配：第一次创建一个临时字符串并绑定到输入参数，然后再把它拷贝进 \texttt{mName}。
\item \texttt{u2}——一次分配：我们以零成本的代价绑定到引用，然后再把它拷贝进成员变量。
\item \texttt{u3}——一次分配：我们以零成本的代价绑定到引用，然后再把它拷贝进成员变量。
\item 你得额外写一个接受右值引用的构造函数，才能在 \texttt{u1} 的情形下省去一次分配，并且这也能在 \texttt{u3} 的情形下省去一次拷贝（因为我们可以从右值引用移动）。
\end{itemize}

对于 \texttt{std::string\_view}：

\begin{itemize}
\item \texttt{u1}——一次分配——输入参数没有拷贝/分配，只有在创建 \texttt{mName} 时才有一次分配。
\item \texttt{u2}——一次分配——为实参创建一个廉价的 \texttt{string\_view}，然后再把它拷贝进成员变量。
\item \texttt{u3}——一次分配——为实参创建一个廉价的 \texttt{string\_view}，然后再把它拷贝进成员变量。
\item 如果你希望在 \texttt{u3} 的情形下省去一次分配，你也得写一个接受右值引用的构造函数，因为你可以从右值引用移动。
\end{itemize}

虽然当你传入字符串字面量时 \texttt{string\_view} 表现更好，但当你针对已有的字符串使用它、或从它移动时，它并没有更好。

然而，自 C++11 引入移动语义以来，通常更好的、也更安全的做法是把字符串按值传入，然后再从中移动。

例如：

\begin{codebox}[cpp]
class UserName {
std::string mName;
public:
UserName(std::string str) : mName(std::move(str)) { }
};
\end{codebox}

现在我们得到如下结果：

对于 \texttt{std::string}：

\begin{itemize}
\item \texttt{u1}——一次分配——为输入实参分配，然后一次移动进入 \texttt{mName}。这比使用 \texttt{const std::string\&} 时我们得到两次内存分配要好，也与 \texttt{string\_view} 的做法类似。
\item \texttt{u2}——一次分配——我们得把值拷贝进实参，然后可以从中移动。
\item \texttt{u3}——没有分配，只有两次移动操作——这比使用 \texttt{string\_view} 和 \texttt{const string\&} 都要好！
\end{itemize}

当你按值传递 \texttt{std::string} 时，代码不仅更简单，而且也不再需要为右值引用编写单独的重载。

完整的代码示例见 \texttt{Chapter string\_view/initializing\_from\_string\_view.cpp}。

这种按值传递的做法与 Scott Meyers 所著《Effective Modern C++》中的第 41 条“对于廉价移动且总是会被拷贝的可拷贝形参，考虑按值传递”是一致的。

然而，\texttt{std::string} 移动起来真的廉价吗？

尽管 C++ 标准并未规定这一点，但通常字符串是通过小字符串优化（SSO）来实现的——字符串对象含有额外的空间，用于在不额外分配内存的情况下容纳字符⁴。这意味着移动一个字符串与拷贝它是一样的。而且由于字符串较短，拷贝也很快。

让我们重新考虑当字符串较短时按值传递的例子：

\begin{codebox}[cpp]
UserName u1{"John"}; // fits in SSO buffer
std::string s2 {"Marc"}; // fits in SSO buffer
UserName u2 { s2 };
std::string s3 {"Marc"}; // fits in SSO buffer
UserName u3 { std::move(s3) };
\end{codebox}

请记住：每一次移动都等同于一次拷贝。

对于 \texttt{const std::string\&}：

\begin{itemize}
\item \texttt{u1}——两次拷贝：一次从输入的字符串字面量拷贝进临时字符串实参，另一次拷贝进成员变量。
\item \texttt{u2}——一次拷贝：已有的字符串被绑定到引用实参，然后我们把它拷贝进成员变量一次。
\item \texttt{u3}——一次拷贝：右值引用以零成本绑定到输入参数，之后我们把它拷贝进成员字段一次。
\end{itemize}

对于 \texttt{std::string\_view}：

\begin{itemize}
\item \texttt{u1}——一次拷贝：输入参数没有拷贝，只有在初始化 \texttt{mName} 时才发生一次拷贝。
\item \texttt{u2}——一次拷贝：输入参数没有拷贝（因为创建 \texttt{string\_view} 很快），然后拷贝进成员变量一次。
\item \texttt{u3}——一次拷贝：\texttt{string\_view} 创建代价很低，从实参到 \texttt{mName} 有一次拷贝。
\item 额外风险：\texttt{string\_view} 可能指向一个已被销毁的字符串。
\end{itemize}

对于 \texttt{std::string}：

\begin{itemize}
\item \texttt{u1}——两次拷贝：输入实参从字符串字面量创建，然后拷贝进 \texttt{mName}。
\item \texttt{u2}——两次拷贝：一次拷贝进实参，然后第二次拷贝进成员。
\item \texttt{u3}——两次拷贝：一次拷贝进实参（移动即意味着拷贝），然后第二次拷贝进成员。
\end{itemize}

如你所见，对于短字符串，按值传递在你传入某个已有字符串时可能会“更慢”一些，仅仅因为你做了两次拷贝而不是一次。另一方面，当编译器看到一个值时，它或许能把代码优化得更好。而且，短字符串拷贝起来很廉价，所以潜在的“变慢”可能根本看不出来。

总而言之，按值传递然后再从字符串实参移动，是更受青睐的方案。你既能得到简洁的代码，又能在较大字符串上获得更好的性能。

一如既往，如果你的代码需要极致的性能，那就必须把所有可能的情形都实测一遍。

\paragraph{其它类型与自动化}

本节讨论的问题也可以推广到其它可拷贝、可移动的类型。如果移动操作很廉价，那么按值传递可能优于按引用传递。

你也可以使用自动化工具，例如 Clang-Tidy，它能够检测出潜在的改进点。Clang Tidy 针对这种用例有一条专门的规则，参见 \texttt{clang-tidy - modernize-pass-by-value}⁵。

\footnotetext{⁴ SSO 并未标准化，且可能发生变化。目前，MSVC（VS 2017）/GCC（8.1）下为 15 个字符，Clang（6.0）下为 22 个字符。对于跨平台代码，依赖基于 SSO 的优化并不可取。}

\footnotetext{⁵ \url{https://clang.llvm.org/extra/clang-tidy/checks/modernize-pass-by-value.html}}

下面是字符串传递以及字符串成员初始化的总结：

\begin{table}[H]
\centering
\begin{tabularx}{\textwidth}{X c c c}
\toprule
传入形参 & \texttt{const string\&} & \texttt{string\_view} & string 传值并移动 \\
\midrule
\texttt{const char*} & 2 次分配 & 1 次分配 & 1 次分配 + 移动 \\
\texttt{const char*}（SSO） & 2 次拷贝 & 1 次拷贝 & 2 次拷贝 \\
左值 & 1 次分配 & 1 次分配 & 1 次分配 + 1 次移动 \\
左值（SSO） & 1 次拷贝 & 1 次拷贝 & 2 次拷贝 \\
右值 & 1 次分配 & 1 次分配 & 2 次移动 \\
右值（SSO） & 1 次拷贝 & 1 次拷贝 & 2 次拷贝 \\
\bottomrule
\end{tabularx}
\end{table}

\section{处理非以空字符结尾的字符串}

如果你从一个字符串得到 \texttt{string\_view}，那么它会指向一段以空字符结尾的内存：

\begin{codebox}[cpp]
std::string s = "Hello World";
std::cout << s.size() << '\n';
std::string_view sv = s;
std::cout << sv.size() << '\n';
\end{codebox}

两条 \texttt{cout} 语句都会打印出 11。

但如果你只持有字符串的一部分呢？

\begin{codebox}[cpp]
std::string s = "Hello World";
std::cout << s.size() << '\n';
std::string_view sv = s;
auto sv2 = sv.substr(0, 5);
std::cout << sv2.data() << '\n'; /// ooops?
\end{codebox}

\texttt{sv2} 应该只包含 \texttt{"Hello"}，但当你访问指向底层内存的指针时，你会拿到指向整个字符串的指针。表达式 \texttt{cout << sv2.data()} 会打印整个字符串，而不仅仅是它的一部分！\texttt{sv2.data()} 返回的是字符串 \texttt{s} 对象内部那个 \texttt{"Hello World"} 字符数组的指针。

当然，当你打印 \texttt{sv2} 时，你会得到正确的结果：

\begin{codebox}[cpp]
std::cout << sv2 << '\n';
// prints "Hello"
\end{codebox}

这是因为 \texttt{std::cout} 能正确处理 \texttt{string\_view} 类型。

这个例子展示了所有假定以空字符结尾的字符串的第三方 API 都可能存在的潜在问题。举几个例子：

\paragraph{使用 printf() 打印}

例如：

\begin{codebox}[cpp]
std::string s = "Hello World";
std::string_view sv = s;
std::string_view sv2 = sv.substr(0, 5);
printf("My String %s", sv2.data()); // oops?
\end{codebox}

相反，你应该使用：

\begin{codebox}[cpp]
printf("%.*s\n", static_cast<int>(sv2.size()), sv2.data());
\end{codebox}

\texttt{.*} 描述的是精度，参见 printf 规范⁶：

\footnotetext{⁶ \url{http://www.cplusplus.com/reference/cstdio/printf/}}

精度并不在格式字符串中指定，而是作为一个额外的整数值实参，放在那个需要被格式化的实参之前。

\paragraph{像 atoi()/atof() 这样的转换函数：}

\begin{codebox}[cpp]
std::string number = "123.456";
std::string_view svNum { number.data(), 3 };
auto f = atof(svNum.data()); // should be 123, but is 123.456!
std::cout << f << '\n';
\end{codebox}

\texttt{atof} 只接受指向以空字符结尾的字符串的指针，因此与 \texttt{string\_view} 不兼容。

要修复这一点，你可以看看 \texttt{from\_chars} 函数（同样是在 C++17 中加入的）：

\begin{codebox}[cpp]
// use from_chars (C++17)
std::string number = "123.456";
std::string_view svNum { number.data(), 3 };
int res = 0;
std::from_chars(svNum.data(), svNum.data()+svNum.size(), res);
std::cout << res << '\n';
\end{codebox}

\paragraph{通用解决方案}

如果你的 API 只支持以空字符结尾的字符串，而你又无法切换到接受额外计数或大小参数的函数，那么你需要把一个视图转换成字符串。

例如：

\begin{codebox}[cpp]
void ConvertAndShow(const char *str)
{
auto f = atof(str);
std::cout << f << '\n';
}
std::string number = "123.456";
std::string_view svNum { number.data(), 3 };
// ... some code
std::string tempStr { svNum.data(), svNum.size() };
ConvertAndShow(tempStr.c_str());
\end{codebox}

\texttt{ConvertAndShow} 只能处理以空字符结尾的字符串，因此我们唯一的办法就是创建一个临时字符串 \texttt{tempStr}，然后再把它传给该函数。

如果你想从 \texttt{string\_view} 创建一个字符串对象，请记得使用 \texttt{.data()} 和 \texttt{.size()}，这样你引用的是底层字符序列中正确的那一段。

可运行示例见：\texttt{Chapter string\_view/string\_view\_null.cpp}。

\section{性能与内存考量}

把 \texttt{string\_view} 加入标准库，其核心思想在于性能与内存消耗。通过借助 \texttt{string\_view}，你可以有效地避免创建大量临时的字符串，从而提升性能。

关于内存：\texttt{string\_view} 通常实现为 \texttt{[ptr, len]}——一个指针以及一个通常用来表示可能大小的 \texttt{size\_t}。

正因如此，你应当把它的大小视为 8 字节或 16 字节（取决于架构是 x86 还是 x64）。

如果我们考虑 \texttt{std::string} 类型，由于常见的小字符串优化，\texttt{std::string} 通常是 24 或 32 字节，也就是 \texttt{string\_view} 的两倍大小。如果字符串比 SSO 缓冲区更长，那么 \texttt{std::string} 会在堆上分配内存。如果 SSO 不被支持（这种情况很少见），那么 \texttt{std::string} 将由指向已分配内存的指针和大小构成。

关于字符串操作的性能。

\texttt{string\_view} 只拥有字符串操作的一个子集，即那些不会修改被引用字符序列的操作。像 \texttt{find()} 这样的函数，应当提供与字符串版本相同的性能。

另一方面，\texttt{substr} 在 \texttt{string\_view} 中只是拷贝两个元素，而字符串则会执行一段内存区间的拷贝。复杂度是 O(1) 对 O(n)。因此，如果你需要切分一个较大的字符串并对这些切片进行处理，\texttt{string\_view} 的实现应当提供更好的速度。

\subsection{常量表达式中的字符串}

\texttt{string\_view} 一个有趣的特性是，它的所有方法都被标记为 constexpr（\texttt{copy}、\texttt{operator<<} 以及为字符串视图特化的 \texttt{std::hash} 函数除外）。借助这一能力，你可以在编译期处理字符串。

例如：

\begin{codebox}[cpp]
// Chapter string_view/string_view_constexpr.cpp
#include <string_view>
int main()
{
using namespace std::literals;
constexpr auto strv = "Hello Programming World"sv;
constexpr auto strvCut = strv.substr("Hello "sv.size());
static_assert(strvCut == "Programming World"sv);
return strvCut.size();
}
\end{codebox}

如果你使用一个现代编译器，例如带了以下选项的 GCC 8.1：\texttt{-std=c++1z -Wall -pedantic -O2}，那么编译出的汇编代码应当形如：

\begin{codebox}
main:
movl
$17, %eax
ret
\end{codebox}

同样一段代码的类似版本，若使用 \texttt{std::string}，则会生成多得多的代码。由于这个例子使用了长字符串，小字符串优化无从施展，于是编译器必须为字符串的内存管理生成 \texttt{new}/\texttt{delete} 的代码。

\section{从 boost::string\_ref 与 boost::string\_view 迁移}

与 C++17 中大多数新类型一样，\texttt{string\_view} 也受到了 Boost 库的启发。Marshall Clow 在 1.53 版（2012 年 2 月）中实现了 \texttt{boost::string\_ref}，随后它演进为 \texttt{boost::string\_view}（加入于 1.61 版——2016 年 5 月）。

\texttt{boost::string\_ref} 与 \texttt{boost::string\_view} 之间的主要区别在于对 constexpr 的支持。

\texttt{boost::string\_view} 实现了与 \texttt{std::string\_view} 相同的功能，并额外增加了几个新函数：

\begin{itemize}
\item \texttt{starts\_with}
\item \texttt{ends\_with}
\end{itemize}

完整的头文件见 \texttt{boost/doc/libs/1\_67\_0/boost/utility/string\_view.hpp}⁷，以及 Boost 的 utility 库⁸。

\footnotetext{⁷ \url{https://www.boost.org/doc/libs/1_67_0/boost/utility/string_view.hpp}}

\footnotetext{⁸ \url{https://www.boost.org/doc/libs/1_67_0/boost/utility/}}

关于 \texttt{string\_ref} 被弃用的讨论链接：“string\_view versus string\_ref”⁹。

\footnotetext{⁹ \url{https://lists.boost.org/Archives/boost/2017/07/236903.php}}

\section{示例}

下面你可以找到两个使用 \texttt{string\_view} 的例子。

\subsection{处理不同的字符串 API}

\texttt{string\_view} 一个有趣的用例，是当你在一段需要与不同字符串实现打交道的代码中使用它时。

例如，你可能拥有来自 MFC 的 \texttt{CString}、来自 C-API 的 \texttt{const char*}、来自 Qt 的 \texttt{QString}，当然还有 \texttt{std::string}。

与其为不同的字符串类型编写重载，不如借助 \texttt{string\_view}！

例如：

\begin{codebox}[cpp]
void Process(std::string_view sv) { }
\end{codebox}

如果你想把 \texttt{Process} 用于不同的字符串实现，那么你要做的只是从你的类型创建一个字符串视图。大多数字符串类型都应当能轻松做到这一点。

例如：

\begin{codebox}[cpp]
// MFC Strings:
CString cstr;
Process(std::string_view{cstr.GetString(), cstr.GetLength()});
// QT Strings:
QString qstr;
Process(std::string_view{qstr.toLatin1().constData()};
// Your implementation:
MySuperString myStr;
// MySuperString::GetData() - returns char*
// MySuperString::Length() - returns length of a string
Process(std::string_view{myStr.GetData(), myStr.Length()});
\end{codebox}

假设性地讲，\texttt{Process()} 本可以实现为 \texttt{Process(const char*, int len)}，但使用 \texttt{string\_view} 后代码更明确、也更简洁。此外，你还拥有 \texttt{string\_view} 的全部可用方法，这样的代码也比 C 风格更方便。

\subsection{字符串分割}

\texttt{string\_view} 可能是字符串分割的一个潜在优化点。如果你拥有一个大而持久的字符串，你或许想创建一个 \texttt{string\_view} 对象的列表，用来映射那个较大字符串中的各个单词。

请注意，这段代码受到了 Marco Arena 的文章——\texttt{string\_view odi et amo}¹⁰ 的启发。

\footnotetext{¹⁰ \url{https://marcoarena.wordpress.com/2017/01/03/string_view-odi-et-amo/}}

\begin{codebox}[cpp]
// Chapter string_view/string_view_split.cpp
std::vector<std::string_view>
splitSV(std::string_view strv, std::string_view delims = " ")
{
std::vector<std::string_view> output;
auto first = strv.begin();
while (first != strv.end())
{
const auto second = std::find_first_of(
first, std::cend(strv),
std::cbegin(delims), std::cend(delims));
if (first != second)
{
output.emplace_back(strv.substr(std::distance(strv.begin(), first),
std::distance(first, second)));
}
if (second == strv.end())
break;
first = std::next(second);
}
return output;
}
\end{codebox}

使用示例：

\begin{codebox}[cpp]
const std::string str {
"Hello
Extra,,, Super, Amazing World"
};
for (const auto& word : splitSV(str, " ,"))
std::cout << word << '\n';
\end{codebox}

程序将输出：

\begin{codebox}
Hello
Extra
Super
Amazing
World
\end{codebox}

该算法遍历输入的 \texttt{string\_view} 并寻找断点——即与分隔符相匹配的字符。然后代码提取该序列的一部分——位于上一次断点与本次新断点之间。这个子视图被存入输出的 vector 中。

关于该实现的一些说明：

\begin{itemize}
\item 该算法的 \texttt{string\_view} 版本假定输入的字符串是持久的，而不是一个临时对象。对于返回的 \texttt{string\_view} vector 要小心，因为它同样指向那个输入字符串。
\item 指令 \texttt{if (first != second)}——可以防止加入空的“单词”，例如在出现多个相邻分隔符（比如双空格）的情况下。
\item 该算法使用了 \texttt{std::find\_first\_of}，但也可以使用 \texttt{string\_view::find\_first\_of}。成员方法不会返回迭代器，而是返回在字符串中的位置。在某些测试中，当分隔符数量较少时，成员方法比 \texttt{std::} 版本更慢¹¹。
\end{itemize}

\footnotetext{¹¹ 相关实验见 C++17 中的 std::string\_view 与 std::string 性能对比，以及 Bartek 的编程博客：C++17 中的 std::string\_view 与 std::string 性能对比，其中有专门的讨论和可能的其他改进。}

\section{小结}

以下是关于 \texttt{std::string\_view} 需要记住的要点：

\begin{itemize}
\item 它是 \texttt{std::basic\_string\_view<charType, traits<charType>>} 的一个特化——其中 \texttt{charType} 等于 \texttt{char}。
\item 它是一个连续字符序列的非拥有视图。
\item 它在末尾可能不包含空终止符。
\item 它可以用来优化代码，并减少对字符串临时拷贝的需求。
\item 它包含了 \texttt{std::string} 中那些不改变底层字符的大部分操作。
\item 它的操作也被标记为 constexpr。
\end{itemize}

但是：

\begin{itemize}
\item 确保底层字符序列仍然存在！
\item 虽然 \texttt{std::string\_view} 看起来像是字符串的一个常量引用，但语言并不会延长那些绑定到 \texttt{std::string\_view} 的返回临时对象的生命期。
\item 当你从 \texttt{string\_view} 构建字符串时，务必记得使用 \texttt{stringView.size()}。\texttt{size()} 方法能正确标记 \texttt{string\_view} 的结尾。
\item 当你把 \texttt{string\_view} 传入接受以空字符结尾的字符串的函数时要小心，除非你确信你的 \texttt{string\_view} 包含了一个空终止符。
\end{itemize}

\subsection{编译器支持}

\begin{table}[H]
\centering
\begin{tabularx}{\textwidth}{X c c c}
\toprule
特性 & GCC & Clang & MSVC \\
\midrule
\texttt{std::string\_view} & 7.1 & 4.0 & VS 2017 15.0 \\
\bottomrule
\end{tabularx}
\end{table}
